% HD-S51. GPT-6 Astra Pro /root/sixfolds; inherited effort not exposed.
% 21 November 2026. Actual kernel, projectors, quotient and polarized corner.
\section{HD-S51: three-dimensional rank-four blocks}
\label{sec:three-block-rank-four-scope}

Let \(K\subset\mathbf C\) have characteristic zero. Let
\(N\subset GL_{3r}(K)\), \(r\ge7\), be finite and absolutely
irreducible, generated by actual matrices of moving rank four,
and containing no nonidentity element of moving rank one, two
or three. Suppose its complex representation has a system of
\(r\) three-dimensional blocks. Then the kernel \(T\) of the
permutation action on these blocks is characteristic in \(N\).
Every original block projector has entries in \(K\), and every
actual \(K\)-linear group normalizing \(N\) permutes those
same descended \(K\)-blocks.

For the natural sixfold problem, apply this to the absolutely
irreducible normal subgroup \(N_4\) generated by all actual
rank-four elements of the original coefficient group. The
sphere budget already gives coefficient rank at least twenty-one,
so all three-dimensional complex block systems have \(r\ge7\).
The construction below gives a literal rank-three quotient source
for the originally specified polarized primitive corner. At period
rank nine its determinant belongs to the already supplied classes:
\([-1]\) or \([-3]\) for \(K=\mathbf Q(i)\), and \([-1]\)
for
\[
 K=\mathbf Q(\sqrt{-d}),\qquad
 d\in\{2,7,11,19,43,67,163\}.                         \tag{S51.1}
\]
The structural assertion applies to every characteristic-zero
\(K\); the displayed determinant conclusion has precisely the
eight-field scope of HD-S40 and HD-G27.

\paragraph{1. The actual local generators and block permutation.}
Write \(V_{\mathbf C}=\bigoplus_{i=1}^rV_i\), with
\(\dim V_i=3\). For a cycle of \(\ell\) blocks, a fixed
vector is determined by its value in the first block and must
be fixed by the return product \(R\). Its exact contribution
to the moving rank is
\[
                 3(\ell-1)+\operatorname{rank}(R-I).
                                                               \tag{S51.2}
\]
Thus an outside rank-four generator exchanges precisely two
blocks. A return defect of rank one would make its square have
moving rank two, with all other blocks fixed, contrary to the
hypothesis. The return product is therefore the identity. The
remaining unit of moving rank occurs on exactly one third block.
If its nonidentity eigenvalue there is \(\lambda\), squaring
gives moving rank one unless \(\lambda^2=1\). Finite order
makes this block action semisimple, and \(\lambda\ne1\), so
\(\lambda=-1\). In suitable ordered coordinates the generator
is exactly
\[
 \left. x\right|_{V_i\oplus V_j}
      =\begin{pmatrix}0&A\\A^{-1}&0\end{pmatrix},
 \qquad \left.x\right|_{V_k}=R_k,
 \qquad R_k^2=I,\quad\operatorname{rank}(R_k-I)=1,
                                                               \tag{S51.3}
\]
with identity on the other blocks and \(k\notin\{i,j\}\).
In particular it is an involution. An inside generator is
nonidentity on at most four blocks. Every generator is therefore
literally the identity on at least one entire block.

The graph of exchanged pairs is connected: otherwise the sum of
the blocks in a connected component is invariant under every
generator, contradicting irreducibility. Its edge transpositions
generate \(S_r\), by successive conjugation along paths. Hence
\[
                              N/T\simeq S_r.              \tag{S51.4}
\]
The third-block involutions in (S51.3) have been retained. No
section of (S51.4), and no simultaneous pure-permutation gauge,
is asserted.

\paragraph{2. The coordinate kernel is irreducible.}
Let \(H_i\) be the stabilizer of \(V_i\) in \(N\).
Its action on \(V_i\) is irreducible. Indeed a proper
invariant subspace \(W\subset V_i\) would give the proper
\(N\)-invariant subspace \(\bigoplus_j n_jW\), where
\(n_jV_i=V_j\); independence of each summand from the choice
of \(n_j\) follows from \(H_i\)-invariance. We have the
actual quotient
\[
                         H_i/T\simeq S_{r-1}.              \tag{S51.5}
\]
Put \(L_i=\rho_i(T)\subset GL(V_i)\). The \(T\)-isotypic
subspaces of \(V_i\) are permuted transitively by \(H_i\):
a sum over a proper orbit would be \(H_i\)-invariant. They
have equal dimensions and equal irreducible multiplicities.
Since \(\dim V_i=3\), a reducible restriction is therefore
either scalar, or the sum of three distinct one-dimensional
characters. In the latter case the action on the three character
lines factors through (S51.5), since \(T\) fixes each line.
A homomorphism from \(S_{r-1}\), \(r-1\ge6\), to \(S_3\)
has image trivial or of order two: its kernel is normal, and
the only other possibility would be an injection of the larger
symmetric group. Neither image is transitive on three labels.
This excludes the nonscalar reducible restriction.

If \(L_i\) is scalar, the projective action of \(H_i\)
factors through \(S_{r-1}\). Its image is trivial, cyclic of
order two, or isomorphic to \(S_{r-1}\). In the first two
cases all its actual matrices are scalars times powers of one
finite-order lift, so they preserve an eigenline, contradicting
irreducibility. The third case cannot be an irreducible finite
projective group in degree three. An imprimitive such group is
monomial on three lines and is solvable. The primitive finite
\(PGL_3(\mathbf C)\) list consists of \(A_5,PSL_2(7),A_6\)
and the solvable Hessian groups of orders at most216. Their
orders are at most360, less than \(|S_{r-1}|\) for \(r\ge7\).
At the endpoint \(r=7\), the listed \(A_6\) is not a
quotient of \(S_6\): such a quotient would have a normal
kernel of order two, whereas the nontrivial proper normal
subgroup of \(S_6\) is \(A_6\). This also deals explicitly
with that order360 entry. Consequently every \(L_i\) is
irreducible.

The exact published projective list used here and below is
Dolgachev--Iskovskikh, \emph{Finite subgroups of the plane
Cremona group}, Theorem4.8. The chapter is published in
\emph{Algebra, Arithmetic, and Geometry}, Progress in
Mathematics269 (2009),443--548, DOI
\url{https://doi.org/10.1007/978-0-8176-4745-2_11}.
The retained author-hosted typeset version, at
\url{https://math.nyu.edu/~tschinke/.manin/final/dolgachev/dolgachev.pdf},
has the complete theorem on its literal printed p.459,
PDF page23. Its pagination starts at437, so this locator is
not identified with the different published pagination. The
original page was freshly visually read. The theorem cites
Blichfeldt's classification; no fresh reconstruction of that
classification proof is claimed here. Its Hessian groups are
\(3^2:SL_2(3),3^2:C_4,3^2:Q_8\); these are solvable since
\(SL_2(3)\) has normal \(Q_8\) and quotient \(C_3\).

\paragraph{3. The full or diagonal projective kernel.}
Choose actual elements \(n_i\in N\) with \(n_iV_1=V_i\)
and use \(n_i|_{V_1}\) to identify the coordinate images
\(L_i\) with a common finite irreducible
\(L\subset GL_3(\mathbf C)\). Then
\[
 T\hookrightarrow L^r
 \quad\hbox{has surjective projection onto each coordinate.}
                                                               \tag{S51.6}
\]
Conjugation by \(N\) permutes these coordinates, with the
actual intervening automorphisms of \(L\). The coordinate
identifications are used only in these group calculations;
no field descent has been assumed.

If \(L\) is solvable, (S51.6) makes \(T\) solvable. A
solvable normal subgroup of \(N\) maps trivially to
\(N/T=S_r\), whose solvable radical is trivial. Hence
\[
                              T=\operatorname{Rad}(N)
                                                               \tag{S51.7}
\]
and \(T\) is characteristic.

Suppose \(L\) is nonsolvable. Its degree-three irreducibility
and the projective list imply
\[
 L/Z\simeq S\in\{A_5,PSL_2(7),A_6\},
       \qquad Z=L\cap\mathbf C^*I_3.                    \tag{S51.8}
\]
Indeed an imprimitive degree-three group is solvable, and
all other primitive projective entries are solvable. Each
listed \(S\) is nonabelian simple and perfect. Projectivize
the actual coordinate maps in (S51.6), putting
\[
 P=\operatorname{im}(T\longrightarrow S^r),\qquad
 Z_0=\ker(T\longrightarrow P)=T\cap Z^r.                \tag{S51.9}
\]
Here \(Z_0\) is abelian, central in \(T\), and normal in
\(N\). The image \(P\) projects onto every factor.

For completeness, every such subgroup of \(S^r\) is a product
of twisted diagonal copies of \(S\) on a partition of the
coordinates. This follows by induction on \(r\). Project
onto the first \(r-1\) factors. The kernel in the last
factor is normal in \(S\), hence is either \(S\) or1.
In the first case the last factor is independent. In the
second, the subgroup is the graph of a surjective homomorphism
from the previous product of diagonal simple factors to \(S\).
The images of those factors are commuting normal subgroups of
\(S\), so exactly one maps isomorphically onto \(S\), and
the others map trivially. This appends the new coordinate to
precisely that diagonal factor, with the displayed isomorphism
as its twist. This is an induction using the actual coordinate
projections and kernels.

Two coordinates belong to the same partition block exactly
when their maps \(P\to S\) have the same kernel. Thus the
partition is preserved by all coordinate permutations and
automorphisms induced by \(N\). Its permutation image is
\(S_r\), so the partition is either all singletons or one
block. Consequently
\[
             P=S^r\quad\hbox{or}\quad
             P\simeq S\text{ is a single twisted diagonal.}
                                                               \tag{S51.10}
\]

\paragraph{4. The diagonal alternative is impossible.}
Assume the second alternative in (S51.10). For any generating
rank-four matrix \(x\), choose a block \(V_j\) on which
\(x\) is literally the identity, as proved in paragraph1.
Then for every \(t\in T\), the \(j\)-th coordinate of
\(xtx^{-1}\) equals that of \(t\). The map from the
single diagonal \(P\) to its \(j\)-th factor is injective.
Thus
\[
                      [x,t]=xtx^{-1}t^{-1}\in Z_0
                              \qquad(t\in T).             \tag{S51.11}
\]
For this fixed \(x\), the map \(t\mapsto[x,t]\) is a
homomorphism \(T\to Z_0\):
\[
 [x,ab]=[x,a]\,a[x,b]a^{-1}=[x,a][x,b],
                                                               \tag{S51.12}
\]
because \(Z_0\subset Z(T)\). It therefore vanishes on
\(T'\). Every generator \(x\) centralizes \(T'\), so
all of \(N\) centralizes it. Absolute irreducibility of
\(N\) makes every such matrix scalar. But the image of
\(T'\) in \(P\) is \(P'=S\), whereas scalar matrices
have trivial image in every projective coordinate. This is a
contradiction. Thus the first alternative holds:
\[
                              T/Z_0\simeq S^r.             \tag{S51.13}
\]

Every solvable normal subgroup of \(N\) first maps trivially
to \(S_r\), and then trivially to the product in (S51.13).
Conversely \(Z_0\) is solvable and normal. Hence
\(\operatorname{Rad}(N)=Z_0\). In \(Q=N/Z_0\), the
subgroup \(B=T/Z_0\) is a minimal nontrivial normal subgroup:
normal subgroups of \(S^r\) are products of its factors,
and the conjugation action on those factors is transitive.
Its centralizer in \(Q\) is trivial. A centralizing element
must have trivial coordinate permutation, since it centralizes
each factor pointwise; it then belongs to \(B\), whose
centre is trivial. Any other minimal normal subgroup of \(Q\)
would have trivial intersection with \(B\), commute with it,
and contradict this centralizer assertion. Thus \(B\) is
the unique minimal nontrivial normal subgroup of \(Q\).
It is characteristic, as is its preimage \(T\) in \(N\).
No splitting of \(Q\to S_r\) is needed or asserted.
Together with (S51.7), this proves characteristicness in every case.

\paragraph{5. The exact projectors descend to \(K\).}
Each original \(V_i\) is an irreducible \(T\)-module by
paragraph2. Equivalence of two such modules is an equivalence
relation on the labels preserved by \(N\), hence by its
full \(S_r\)-permutation action. Two-transitivity implies
that they are either all equivalent or pairwise inequivalent.

In the first case choose actual \(T\)-intertwiners
\(f_i:V_1\to V_i\). The map
\[
 V_1\otimes\mathbf C^r\longrightarrow V_{\mathbf C},
       w\otimes e_i\longmapsto f_i(w)
                                                               \tag{S51.14}
\]
is an isomorphism with inverse given by the coordinate
projections followed by \(f_i^{-1}\). It identifies the
image algebra of \(T\) with \(\operatorname{End}(V_1)\otimes I_r\).
Every \(n\in N\) normalizes this algebra. Take the
three-dimensional-factor matrices of its transported matrix
units, choose a nonzero vector in the first image line and
apply the other image matrix units;
this constructs an inner intertwining matrix on \(V_1\).
After removing it the remaining matrix lies in the commutant,
namely \(I_3\otimes\operatorname{End}(\mathbf C^r)\).
Thus (S51.14) is an actual tensor structure preserved by \(N\).
HD-R17 excludes this for an absolutely irreducible rank-four-generated
group in dimension \(3r>16\). Therefore the modules are
pairwise inequivalent.

Let \(\chi_i\) be the actual character on \(V_i\). The
matrices
\[
 p_i=\frac3{|T|}\sum_{t\in T}\chi_i(t^{-1})\rho(t)
                                                               \tag{S51.15}
\]
are exactly the original block projectors. Character averaging
on \(\operatorname{Hom}(V_i,V_j)\) gives restriction
\(\delta_{ij}I_{V_j}\). Averaging matrix coefficients gives
all matrix units supported on an individual \(V_i\), so the
literal restriction map identifies
\[
 \mathbf C\rho(T)=\bigoplus_{i=1}^r\operatorname{End}(V_i).
                                                               \tag{S51.16}
\]
Thus (S51.15) is precisely the complete primitive central
idempotent set of the original image algebra.

Take \(F=K(\zeta_{\exp T})\). All matrices (S51.15) lie
over \(F\). Every \(\gamma\in\operatorname{Gal}(F/K)\)
permutes this exact idempotent set and commutes with the actual
conjugation by \(N\), whose matrices are over \(K\). The
centralizer of the natural \(S_r\) in the permutation group
of its \(r\) labels is trivial: a commuting permutation
sends a chosen label to a label fixed by its full point
stabilizer, which fixes only that label. Hence every \(\gamma\)
fixes every \(p_i\), and
\[
                  p_i\in M_{3r}(K),\qquad \dim_Kp_iV=3.
                                                               \tag{S51.17}
\]
As \(T\) is characteristic in \(N\), any actual overgroup
normalizing \(N\) normalizes its image algebra and permutes
these same central projectors. Their images are orthogonal for
the original invariant positive coefficient form, since an
off-diagonal pairing would intertwine inequivalent \(T\)-modules.
This proves the structural assertion.

\paragraph{6. The actual smaller cover and the original corner.}
Now let \(K\) be one of the eight imaginary quadratic fields
specified in (S51.1), including \(\mathbf Q(i)\). Return to
the natural-source setting, with the original finite
\(G\subset GL_K(V)\), \(m=3r\), exact character field
\(K\), and positive invariant coefficient Hermitian form
\(h\). Put \(U=p_1V\), \(H=\operatorname{Stab}_G(U)\),
\(\sigma=\rho|_H\) on \(U\), and
\(H_0=\ker\sigma\). The actual induction map and inverse are
\[
 F_v(x)=p_1\rho(x)v,\quad F(hx)=\sigma(h)F(x),\qquad
 F\longmapsto\sum_{x\in H\backslash G}\rho(x)^{-1}F(x).
                                                               \tag{S51.18}
\]
The inverse is independent of coset representatives. The two
composites are identities because the conjugates of \(p_1\)
are orthogonal projections summing to \(I\); equivalently the
sheet labelled \(Ht_j^{-1}\) corresponds to \(t_jU\).
Right monodromy on the functions is \(F(x)\mapsto F(xg)\).
The stabilizer representation is absolutely irreducible: a
proper invariant subspace would induce a proper one of the
original module. It has exact field \(K\), since induced
traces give
\(K\subset\mathbf Q(\chi_\sigma(H))\subset K\).

Let \(\widetilde C\) be the original full-\(G\) curve and take
\[
 Y=\widetilde C/H,\qquad C'=\widetilde C/H_0,
             \qquad \alpha:\widetilde C\longrightarrow C'.
                                                               \tag{S51.19}
\]
Then \(C'\to Y\) is the actual full-\(H/H_0\) cover.
On a cycle of quotient sheets, the fixed-vector equations in
(S51.18) are exactly the fixed-vector equations for the return
monodromy on \(U\). Thus the induction identifies invariant
stalks at every branch point, the middle extensions and their
first parabolic cohomology. It identifies their Hodge structures
through these actual finite curve maps.

For \(v\in U(K)\setminus\{0\}\), put \(q=h(v,v)\).
The actual rational corners agree inside \(\mathbf Q[G]\):
\[
 \begin{aligned}
 e_G&=\frac{2m}{|G|q}\sum_{g\in G}
       \operatorname{Re}h(v,\rho(g^{-1})v)g\\
    &=\frac6{|H|q}\sum_{a\in H}
       \operatorname{Re}h(v,\sigma(a^{-1})v)a=e_H.
 \end{aligned}                                               \tag{S51.20}
\]
Indeed the coefficients off \(H\) vanish by block
orthogonality, and \([G:H]=r\). The coefficients of \(e_H\)
are constant on \(H_0\)-cosets, so it is the averaged
inflation of the corresponding rational corner of \(H/H_0\).
The complex Schur pairing factor for the coefficient map
normalized by \(q\) is \(|G|/(mq)\); for its actual real
map obtained by adding the conjugate it is
\[
                       \frac{2|G|}{mq}=\frac{2|H|}{3q}.
                                                               \tag{S51.21}
\]
These equalities specify the same positive coefficient
normalization on both sides of (S51.18).

On Jacobians, pullback and norm give the exact identities
\[
 \alpha_*\alpha^*=[|H_0|],\qquad
        (\alpha^*)^*\theta_{J\widetilde C}
                            =|H_0|\theta_{JC'}.
                                                               \tag{S51.22}
\]
After integer clearing of the rational projectors,
\(\alpha^*\) restricts to a surjective isogeny from the
connected primitive corner of \(JC'\) onto the corner
(S51.20) in \(J\widetilde C\). Its kernel is finite by
the first identity, and equality of the images follows from
the averaged corner identity and the parabolic comparison.

For an initially prescribed line \(Kw\subset V\), including
one outside \(U\), put \(q_w=h(w,w)\) and
\[
 A(x)=v\,h(w,x)/q_w,\qquad
 b=\frac{2m}{|G|}\sum_{g\in G}
       \operatorname{Re}\operatorname{tr}(A\rho(g^{-1}))g.
                                                               \tag{S51.23}
\]
Schur averaging gives the literal operator \(A\) on the
original coefficient block, its conjugate on the conjugate
block, and zero on every other block. Thus
\[
                         b^\dagger b=(q/q_w)e_w.           \tag{S51.24}
\]
For an integer \(M\) clearing its group-algebra denominator,
the actual endomorphism \(Mb\) restricts to an isogeny of
the original primitive factor onto the corner in (S51.20),
with polarization multiplier \(M^2q/q_w\). All maps are
\(K\)-linear on the specified factors, with the fixed
Jacobian action convention of HD-S23. The combined positive
rational similitudes multiply a six-dimensional Hermitian
determinant by a sixth power, which equals the norm of the
rational cube. They preserve its ordinary determinant class.
The original action, polarization comparison and corner have
therefore all been specified.

\paragraph{7. The rank-three source endpoint.}
The parabolic dimension on \(Y\) is six by (S51.18).
With \(g=g(Y)\), omit identity coefficient inertia and write
\(b_Y\) for the remaining branch count. Absolute irreducibility
kills the invariant and coinvariant terms, so
\[
               6=3(2g-2)+\sum_{j=1}^{b_Y}c_j,
                             \qquad c_j\ge1.               \tag{S51.25}
\]
For \(g=2\), the cover system is unramified and the base
moduli dimension is three. For \(g=1\), \(b_Y\le6\) and
the marked-base dimension is at most six. Genus at least three
is impossible. For \(g=0\), \(\sum c_j=12\), so
\(b_Y\le12\) and the marked-base dimension is at most nine.
Finite-monodromy systems have discrete choices over a fixed
pointed base. The projector and isogeny matrices above are
flat after fixing the finite cover markings. Therefore the
original period differential factors through this actual
rank-three quotient family.

Period rank nine consequently forces \(Y=\mathbf P^1\),
\(b_Y=12\), and \(c_j=1\) for every branch entry. The
actual full coefficient group \(H/H_0\) is generated by
these twelve reflections. This is a rank-three instance of
the complete low-rank natural-source conclusions of HD-S40,
with its seven-field extension and HD-G27. They give precisely
the already supplied determinant alternatives stated in
(S51.1); the explicitly mapped original factor has the same
class by (S51.22)--(S51.24). This closes the three-dimensional
complex-block branch of the absolutely irreducible \(N_4\)
source without replacing any of its actual groups or corners.
